\documentclass[11pt,a4paper]{article}
\usepackage[margin=1in]{geometry}
\usepackage{amsmath,amssymb,amsthm}
\usepackage{booktabs}
\usepackage{hyperref}
\hypersetup{colorlinks=true,linkcolor=blue,urlcolor=blue,citecolor=blue}

\title{\textbf{FermatQuinticJacobian: Lean 4 Certificate for the Jacobian Smoothness of the Fermat Quintic Fourfold}}
\author{Benjamin Stanley Frohman\\
\small ORCID: \href{https://orcid.org/0009-0006-7068-3718}{0009-0006-7068-3718} \quad|\quad
GitHub: \href{https://github.com/BenFrohman/FermatQuinticJacobian}{\texttt{BenFrohman/FermatQuinticJacobian}}\\
\small Dedicated 1-to-1 Zenodo DOI: \href{https://doi.org/10.5281/zenodo.23652994}{\texttt{10.5281/zenodo.23652994}}}
\date{Earliest Commit Provenance: 2026-10-03 12:49:57 PM CDT (UTC-5)\\
Siloed Archival Edition: October 25, 2026 (America/Chicago CDT)}

\usepackage{longtable}
\usepackage{url}
\begin{document}
\maketitle

\begin{abstract}
Standalone Lean 4 formalization (`FermatQuinticJacobian.lean`) certifying the non-degeneracy of the partial derivatives and Jacobian smoothness criterion for the Fermat quintic fourfold V(x\_0\textasciicircum{}5 + ... + x\_5\textasciicircum{}5) subset P\textasciicircum{}5.
\end{abstract}

\section{Scope, Scientific Honesty, and Non-Claim Boundary}
Certifies Jacobian smoothness of the Fermat quintic host variety in Lean 4. Explicitly states that Jacobian smoothness is a prerequisite geometric property and not a resolution of the Hodge Conjecture.

\begin{itemize}
  \item \textbf{Repository:} \texttt{https://github.com/BenFrohman/FermatQuinticJacobian}
  \item \textbf{Dedicated 1-to-1 Zenodo DOI:} \texttt{10.5281/zenodo.23652994}
  \item \textbf{Earliest Git Commit SHA:} \texttt{777a2e97451ec4a366c03bbacb4333c001bb833d} (2026-10-03 12:49:57 PM CDT (UTC-5))
  \item \textbf{License \& Copyright:} Copyright \textcopyright\ 2026 Benjamin Stanley Frohman. Licensed under the Apache License, Version 2.0.
\end{itemize}

\section{Technical Overview \& Repository Specification}
Copyright (c) 2026 Benjamin Stanley Frohman (`@BenFrohman`).

Released under the Apache License, Version 2.0. See `LICENSE` and `NOTICE`.

A Lean 4 certificate, checked against Mathlib `v4.22.0`, that the Fermat quintic

```text

G = X0\textasciicircum{}5 + X1\textasciicircum{}5 + X2\textasciicircum{}5 + X3\textasciicircum{}5 + X4\textasciicircum{}5 + X5\textasciicircum{}5

```

has partial derivatives

```text

pderiv i G = 5 * X i \textasciicircum{} 4

```

Those forms vanish together only at the origin of affine 6-space. By the Jacobian

criterion, `V(G)` in `P\textasciicircum{}5` is smooth.

The theorem in the repository is `Hodge.FermatQuintic.partial\_eq`.

This is not a proof of the Hodge conjecture. It does not construct a cycle from a

Hodge class. It does not show that `H\textasciicircum{}\{3,1\}` vanishes: on this fourfold that

primitive piece is the degree-4 part of `C[X\_i]/(X\_i\textasciicircum{}4)`, of dimension 120.

The non-homogeneous expression

`sum x\_i\textasciicircum{}5 - 5 x0 x1 x2 x3 x4 x5` is not used. The product of six variables has

\section{Cryptographic Provenance \& File Inventory}
Every source artifact in \texttt{BenFrohman/FermatQuinticJacobian} is audited via \texttt{titan-audit}
(\texttt{audit/TITAN\_AUDIT\_CERTIFICATE.json}) and sealed under a deterministic SHA-256 Merkle tree via
\texttt{sealguard} (\texttt{PROVENANCE\_MERKLE\_ROOT.json}) signed with Ed25519 key fingerprint
\texttt{b428fb72676f1c52}.

\begin{center}
\small
\begin{tabular}{p{0.52\textwidth}rr}
\toprule
\textbf{Repository Path} & \textbf{Size} & \textbf{SHA-256 Prefix} \\
\midrule
\texttt{FermatQuinticJacobian.lean} & 1871 B & \texttt{d94c40a67456932d...} \\
\texttt{LICENSE} & 11355 B & \texttt{9da72c2bee4f20ba...} \\
\texttt{NOTICE} & 638 B & \texttt{1d6c88adb668f721...} \\
\texttt{README.md} & 1106 B & \texttt{775cf9e9962b7aa2...} \\
\texttt{lakefile.toml} & 288 B & \texttt{3f8f4cf04ccf5e35...} \\
\texttt{lean-toolchain} & 25 B & \texttt{8f6680f8389e8ade...} \\
\bottomrule
\end{tabular}
\end{center}

\vfill
\noindent\rule{\textwidth}{0.4pt}\\
\footnotesize
\textbf{Attribution \& Licensing Notice:} Copyright \textcopyright\ 2026 Benjamin Stanley Frohman (ORCID: \texttt{0009-0006-7068-3718}).
Licensed under the Apache License, Version 2.0.
\appendix
\section{Verification certificate}\label{app:verification-certificate}

The formal development was machine-checked as follows.

{\raggedright
\begin{itemize}
\item Repository:\\ {\footnotesize\url{https://github.com/BenFrohman/FermatQuinticJacobian}}
\item Verified commit:\\ {\footnotesize\texttt{8ca201272348ac9047601ea60e4567d3e2c984a9}}
\item Lean toolchain: \texttt{leanprover/lean4:v4.22.0}
\item Mathlib: tag \texttt{v4.22.0}, commit (from \texttt{lake-manifest.json})\\ {\footnotesize\texttt{79e94a093aff4a60fb1b1f92d9681e407124c2ca}}
\item Continuous integration run (build, axiom audit, \texttt{sorry}/\texttt{axiom} scan; passed):\\ {\footnotesize\url{https://github.com/BenFrohman/FermatQuinticJacobian/actions/runs/41568055306}}
\item SHA-256 of the GitHub source tarball of the verified commit, downloaded twice with identical result:\\ {\footnotesize\url{https://github.com/BenFrohman/FermatQuinticJacobian/archive/8ca201272348ac9047601ea60e4567d3e2c984a9.tar.gz}}\\
{\footnotesize\texttt{c03d038daf33886f9f0e87cd8d9e2a49e937ff88ba91dc175be345e6939c0aa6}}
\end{itemize}\par}

The tarball hash refers to the verified commit rather than to a release tag, since a document cannot contain the hash of an archive that includes itself. The Lean sources are unchanged between that commit and the release tag.

\textbf{The development contains no \texttt{sorry} and no user-declared axioms.} Every one of the 3 theorems and lemmas listed below depends only on Lean's standard axioms \texttt{propext}, \texttt{Classical.choice} and \texttt{Quot.sound} (or on none). The table is the verbatim output of \texttt{lake env lean AxiomAudit.lean} (\texttt{\#print axioms}) on the verified commit.

{\footnotesize
\begin{longtable}{p{0.55\linewidth}p{0.38\linewidth}}
\hline
Lean declaration & \texttt{\#print axioms} \\
\hline\endhead
\path{Hodge.FermatQuintic.partial_eq} & \texttt{[propext, Classical.choice, Quot.sound]} \\
\path{Hodge.FermatQuintic.partial_off} & \texttt{[propext, Classical.choice, Quot.sound]} \\
\path{Hodge.FermatQuintic.partial_on} & \texttt{[propext, Classical.choice, Quot.sound]} \\
\hline
\end{longtable}}

\end{document}
